\section{Background and Key Concepts}
\label{sec:background}

\subsection{PINN loss structure}
\label{sec:bg-loss}
We consider the diffusion equation
\begin{equation}
  \label{eq:pde}
  \frac{\partial T}{\partial t} \;=\; \kz\,\frac{\partial^2 T}{\partial z^2},
  \qquad z\in[0,2],\ t\in[0,5],
\end{equation}
with diffusivity $\kz=0.05$ in forward problems and $\kz$ estimated from data in inverse problems.
Equation~\eqref{eq:pde} says how fast the temperature changes in \emph{time} (left) equals the diffusivity
times its curvature in \emph{depth} (right); checking whether a network's prediction obeys it means computing
the derivatives each side needs, which is what Section~\ref{sec:bg-stages} times.
Let $\Th(z,t)$ be the network output with trainable parameters $\theta$ and input $x=(z,t)$.
The PINN loss combines four mean-squared terms,
\begin{equation}
  \label{eq:loss}
  \Lloss(\theta) \;=\; \Lloss_{\text{PDE}} + \Lloss_{\text{BC}} + \Lloss_{\text{IC}} + \Lloss_{\text{data}},
\end{equation}
namely the PDE residual at interior collocation points, the boundary-condition residual, the misfit to the initial
condition, and, for inverse problems only, the misfit to observations; all terms have unit weight.

\paragraph{Initial and boundary conditions.}
Closed-form solutions exist for all three boundary-condition types. Table~\ref{tab:bc} lists the conditions
imposed. The initial condition is the exact solution at $t=0$; for the Dirichlet problem the exact solution
vanishes at both boundaries, so $\Lloss_{\text{BC}}$ penalises $\Th(0,t)$ and $\Th(2,t)$ directly. The Neumann and
Robin conditions involve $\partial\Th/\partial z$ at the boundaries and therefore already require a first-order
input gradient (Section~\ref{sec:bg-first}).

\begin{table}[h]
\centering
\caption{Initial and boundary conditions of the three benchmark problems. Each boundary condition holds at
$z=0$ and $z=2$ for all $t\in[0,5]$.}
\label{tab:bc}
\small
\begin{tabular}{@{}lll@{}}
\toprule
Type & Boundary condition & Initial condition $T(z,0)$ \\
\midrule
Dirichlet & $T=0$ & $\sin\frac{\pi z}{2}+\sin\pi z$ \\
Neumann   & $\partial T/\partial z=0$ & $\cos\frac{\pi z}{2}+\cos\pi z$ \\
Robin     & $\partial T/\partial z+T=0$ &
  $e^{-z}+\frac{\pi}{2}\cos\frac{\pi z}{2}-\sin\frac{\pi z}{2}+\pi\cos\pi z-\sin\pi z$ \\
\bottomrule
\end{tabular}
\end{table}

\paragraph{Training points.}
We sample $2540$ interior collocation points uniformly in $[0,2]\times[0,5]$, $40$ points on each of the two
boundaries, and $160$ points at $t=0$. Inverse problems additionally use $12$ observations on a uniform
$4\times3$ grid in $(z,t)$.

\subsection{Four cumulative stages}
\label{sec:bg-stages}
Training and evaluating a PINN involves four computations of increasing cost, each reusing the computation
graph of the one before it. We name them here as they appear in the timing results (Section~\ref{sec:setup})
and summarise their role before deriving each in detail below.
\begin{description}
  \item[forward] evaluates the network once, $\Th(z,t)$, with no derivatives at all -- the cost of using the
  trained network to make a single prediction.
  \item[first\_grad] additionally computes the \emph{first-order input gradient}
  $\nabla_x\Th=(\partial\Th/\partial z,\ \partial\Th/\partial t)$ (Section~\ref{sec:bg-first}), needed for the
  time-derivative term on the left of Eq.~\eqref{eq:pde} and for the Neumann and Robin boundary residuals.
  \item[second\_grad] additionally computes the \emph{second-order input gradient} $\partial^2\Th/\partial z^2$
  (Section~\ref{sec:bg-second}), the diffusion term on the right of Eq.~\eqref{eq:pde}. Only once this is
  available can the PDE residual $r_\theta$ (below) be evaluated at all, so second\_grad is the cost of
  \emph{checking} whether a prediction satisfies the governing equation.
  \item[param\_grad] additionally computes the \emph{parameter gradient} $\partial\Lloss_{\text{PDE}}/\partial\theta$
  (Section~\ref{sec:bg-param}), obtained by treating $r_\theta$ as part of the loss and backpropagating once
  more, to the network's weights rather than to its input. Where the first three stages only check the
  network, param\_grad is what the optimiser actually runs at every training step to \emph{correct} it.
\end{description}
Because the stages are cumulative, second\_grad already includes the work of forward and first\_grad; a ratio
such as the second\_grad-to-forward ratio is therefore the cost of \emph{everything} the PDE residual needs,
relative to a single plain prediction, not the cost of the second derivative in isolation.

\subsection{First-order input gradient}
\label{sec:bg-first}
The PDE residual is
\begin{equation}
  r_\theta(z,t) \;=\; \frac{\partial \Th}{\partial t} \;-\; \kz\,\frac{\partial^2 \Th}{\partial z^2}.
\end{equation}
Evaluating it requires the \emph{first-order gradient} of the output with respect to its input,
$\nabla_x \Th = (\partial\Th/\partial z,\ \partial\Th/\partial t)$, which PINNs obtain by
automatic differentiation of the network~\cite{raissi2019,lu2021deepxde,baydin2017}. It supplies the time derivative in $r_\theta$
and the boundary residuals of the Neumann ($\partial\Th/\partial z = 0$) and Robin
($\partial\Th/\partial z + \Th = 0$) problems. This is the \texttt{first\_grad} stage of
Section~\ref{sec:bg-stages}. It must be distinguished from the \emph{parameter gradient}
$\partial\Lloss/\partial\theta$ used by the optimiser (Section~\ref{sec:bg-param}), which is a further
backward pass through the graph of the input gradients.

\subsection{Second-order input gradient}
\label{sec:bg-second}
The diffusion term requires $\partial^2\Th/\partial z^2$, obtained by differentiating
$\partial\Th/\partial z$ once more with respect to $z$. We compute both derivatives by reverse-mode
automatic differentiation~\cite{baydin2017} in PyTorch~\cite{paszke2017autodiff}, retaining the computation
graph after the first pass so that it can be differentiated again; differentiating through a gradient
with respect to the input in this way is known as double backpropagation~\cite{drucker1992}. The second-order gradient therefore costs a second backward pass through the
graph created by the first, on top of the forward pass and the first-order pass. This defines the
\texttt{second\_grad} stage of Section~\ref{sec:bg-stages}: once it is available, $r_\theta$ can be formed and
squared into $\Lloss_{\text{PDE}}$.

\subsection{Parameter gradient}
\label{sec:bg-param}
Neither $\nabla_x\Th$ nor $\partial^2\Th/\partial z^2$ is what the optimiser uses to update the network: it needs
the gradient of the loss with respect to the network's own trainable parameters, $\partial\Lloss/\partial\theta$.
Restricting to the PDE term, $\Lloss_{\text{PDE}} = \tfrac1N\sum_i r_\theta(z_i,t_i)^2$, computing
$\partial\Lloss_{\text{PDE}}/\partial\theta$ means backpropagating through $r_\theta$ itself -- through the
graph already built for the first- and second-order input gradients -- one more time, to $\theta$ rather than
to $x$. This is the \texttt{param\_grad} stage: approximately the cost of the PDE term in one optimiser step,
and not to be confused with the \emph{input} gradients of Sections~\ref{sec:bg-first}--\ref{sec:bg-second},
which check the network rather than update it. The boundary, initial-condition and (for inverse problems)
data terms of $\Lloss$ in Eq.~\eqref{eq:loss} contribute further, cheaper backward passes of their own that we
do not separately time, so a full training step costs somewhat more than param\_grad.

\subsection{Variational quantum circuits}
\label{sec:bg-vqc}
A VQC acts on $n$ qubits and consists of an encoding of classical features into gate angles followed by
$L$ layers of parametrised single-qubit rotations and entangling gates; its outputs are expectation values
of observables~\cite{benedetti2019,cerezo2021}. The trainable quantum parameters are the rotation angles;
$n$ is the \emph{width} and $L$ the \emph{depth} of the circuit. The encoding gates determine which functions
of the input the circuit can represent~\cite{schuld2021encoding}; we use angle embedding
(Section~\ref{sec:models}). VQCs are the model class of the noisy intermediate-scale quantum (NISQ)
era~\cite{preskill2018}, but in practice they are mostly run on classical simulators, and so are ours; we do
not use quantum hardware.
A statevector simulator stores $2^n$ complex amplitudes and, in the worst case, applies each gate in
$O(2^n)$ operations, so memory and time grow exponentially with $n$ and linearly with the number of gates,
hence with the depth $L$~\cite{xu2023}.

\subsection{How circuit gradients are obtained}
\label{sec:bg-grads}
On a classical simulator, as in our implementation, gradients are obtained by backpropagation: each backward
pass traverses the operations that simulate the circuit, so the cost of every stage grows with the simulated
circuit. On quantum hardware this is not possible, because intermediate quantum states cannot be measured and
stored for reuse without destroying them~\cite{schuld2019}. Gradients of circuit outputs are instead obtained
with the parameter-shift rule~\cite{mitarai2018,schuld2019}: for a gate generated by a Pauli operator, the
derivative with respect to its angle is computed from two evaluations of the circuit with that angle shifted.
Higher-order derivatives follow by iterating the rule, with up to $2^d$ shifted evaluations per derivative
entry of order $d$ (four for the Hessian)~\cite{mari2021}. The same rule applies to the encoding gates,
which gives derivatives of a quantum model with respect to its input~\cite{kyriienko2021}. In our hybrid network
the circuit input is produced by classical layers, so input derivatives of the full network combine the chain
rule through those layers with derivatives of the circuit.
We measure the simulator case only and make no claim that simulator ratios carry over to hardware.

\subsection{Notation}
\begin{table}[h]
\centering
\caption{Notation and terminology.}
\label{tab:notation}
\small
\begin{tabular}{@{}p{0.24\linewidth} p{0.68\linewidth}@{}}
\toprule
Symbol / term & Meaning \\
\midrule
$x=(z,t)$ & network input: depth and time \\
$\Th$ & network output, temperature, with parameters $\theta$ \\
$\kz$ & thermal diffusivity \\
$\nabla_x\Th$ & first-order input gradient (\texttt{first\_grad} stage) \\
$\partial^2\Th/\partial z^2$ & second-order input gradient (\texttt{second\_grad} stage) \\
$\partial\Lloss_{\text{PDE}}/\partial\theta$ & parameter gradient (\texttt{param\_grad} stage; the optimiser's own backward pass) \\
$n$, $L$ & number of qubits (width), number of circuit layers (depth) \\
forward, first\_grad, second\_grad, param\_grad & the four cumulative timed stages (Sections~\ref{sec:bg-stages}, \ref{sec:setup}) \\
\bottomrule
\end{tabular}
\end{table}
